% ══════════════════════════════════════════════
% config.tex — Configuración compartida
% ══════════════════════════════════════════════

\usepackage[utf8]{inputenc}
% ── Idioma ──
% Requiere texlive-langspanish (spanish.ldf + hyph-es.tex). Sin ese paquete
% \usepackage[spanish,es-nodecimaldot]{babel} aborta con "Unknown option 'spanish'".
% babel-espanol mete un espacio fino antes de '%' incluso en modo math, lo que rompe
\usepackage[spanish,es-nodecimaldot]{babel}
% babel-espanol mete un espacio fino antes de '%' incluso en modo
% math, lo que rompe \,74\,\% con "Incompatible glue units".
\spanishplainpercent
\usepackage{csquotes}
\usepackage[T1]{fontenc}
\usepackage[left=6cm,right=1.5cm,top=2.5cm,bottom=2.5cm]{geometry}
\usepackage[colorlinks=true, linkcolor=black, citecolor=black, urlcolor=blue]{hyperref}
\usepackage{amsmath}
\usepackage{amssymb}
\usepackage{wasysym}
\usepackage{booktabs}
\usepackage[version=4]{mhchem}
\usepackage{graphicx}
\usepackage{xcolor}
\usepackage{float}
\usepackage{enumitem}
\usepackage[export]{adjustbox}
\usepackage{epigraph}
\usepackage{tikz}
% Imprescindible con babel español: los caracteres < y > se vuelven activos y
% rompen la sintaxis de claves de TikZ (p. ej. ">=stealth"), provocando
% "Argument of \language@active@arg> has an extra }".
\usetikzlibrary{babel}
\usepackage{helvet}
\usepackage{mathptmx}
\usepackage{array}
\usepackage[table]{xcolor}
\usepackage{caption}
\usepackage{mathtools}

% ── Encabezados y pies de página ──
\usepackage{fancyhdr}

% ── Marca de agua vertical en el margen ──
\usepackage{draftwatermark}
\SetWatermarkColor{blue}
\SetWatermarkAngle{90}
\SetWatermarkFontSize{9pt}
\SetWatermarkHorCenter{0.03\paperwidth}
\SetWatermarkText{Universidad Nacional Aut\'{o}noma de M\'{e}xico, Facultad de Ingenier\'{i}a}

% ── Bibliografía ──
\usepackage[
  backend=biber,
  style=authoryear,
  sorting=nyt,
  maxcitenames=2,
  uniquelist=false,
  hyperref=true
]{biblatex}
\addbibresource{references.bib}

% ── Unificar hipervínculos en citas (author+year en un solo link) ──
\DeclareCiteCommand{\parencite}[\mkbibparens]
  {\usebibmacro{prenote}}
  {\usebibmacro{citeindex}%
   \printtext[bibhyperref]{\usebibmacro{cite}}}
  {\multicitedelim}
  {\usebibmacro{postnote}}

\DeclareCiteCommand{\textcite}
  {\usebibmacro{prenote}}
  {\usebibmacro{citeindex}%
   \printtext[bibhyperref]{\usebibmacro{textcite}}}
  {\multicitedelim}
  {\usebibmacro{textcite:postnote}}

% ══════════════════════════════════════════════
%  NOTAS DE COMPILACIÓN (leer antes de tocar)
% ══════════════════════════════════════════════
%
% Dependencias del sistema (Arch):
%   sudo pacman -S texlive-langspanish biber
%
% 1) Backend: se usa biber, NO bibtex. Con backend=bibtex, biblatex 3.21 emite
%    "Using fall-back bibtex backend" en cada compilación y ese aviso no se
%    puede silenciar desde el preámbulo. Migrado a biber el 2026-10-05.
%
% 2) biber NO está en /usr/bin en Arch: el paquete lo instala como
%    /usr/bin/vendor_perl/biber. Si latexmk se queja de que no encuentra
%    biber, exporta el path:
%      export PATH="$PATH:/usr/bin/vendor_perl"
%    (ya está añadido a ~/.bashrc).
%
% 3) latexmk detecta biber solo, así que el comando normal es:
%      latexmk -pdf main.tex
%    Con -bibtex también funciona, pero es innecesario y confuso: el flag
%    venía de cuando el backend era bibtex y latexmk -pdf a secas NO lo
%    ejecutaba, dejando las citas indefinidas sin avisar.
%
% 4) Tipos de entrada: usa @online para páginas web. @webpage no existe en
%    biblatex y produce "No driver for 'webpage'" en cada compilación.
%    Los .bib de la vault se migraron de @webpage a @online.
